\section{Related Work}
\subsection{Visual and Model-Driven Infrastructure}
Winery implements a topology editor, reusable component definitions, and a repository for TOSCA descriptions~\cite{winery2013}. Its demonstration establishes the feasibility of visual typed composition, while deployment execution is delegated to other components. OpenTOSCA supplies a runtime for packaged application topologies, implementation artifacts, and management plans~\cite{opentosca2013}. These systems are relevant because they separate the representation a user edits from the mechanism that executes it. They prevent a broad novelty claim for CloudCanvas's canvas or deployment graph. Their demonstration evidence also should not be interpreted as a latency baseline for an AWS SDK application developed in a different environment.

CAMEL extends modeling beyond topology to requirements, providers, metrics, scaling, and runtime instances~\cite{camel2019}. Its practitioner study considers perceived usability, rather than proving a universal deployment-speed advantage. The Essential Deployment Metamodel extracts common concepts from deployment technologies and explains limits in mapping among them~\cite{edmm2020}. This highlights a central restriction of CloudCanvas: resource relationships have application-specific meaning, and a visual line is not automatically portable across tools. DOML further separates application, abstract-infrastructure, and concrete-infrastructure models, generates IaC, and uses a Z3-backed checker for model properties~\cite{doml2024}. CloudCanvas's schema and graph checks are narrower than this property-verification approach. Its smaller representation is a design choice whose usability consequences require evidence.

\subsection{AI Generation, Intent, and Deployability}
IaC-Eval provides 458 AWS/Terraform scenarios paired with intent specifications, using a speculative deployment graph and OPA checks to distinguish syntactic acceptance from requested behavior~\cite{iaceval2024}. CloudCanvas currently lacks an equivalent independent intent oracle. Multi-IaC-Eval examines changes to existing templates across multiple formats, combining static tools and semantic review~\cite{multi2025}. Its exclusion of live deployment testing is material: high static compliance cannot be relabeled provider success. It also motivates preservation checks for future graph editing, because CloudCanvas's current AI application action replaces a complete graph rather than performing a verified incremental mutation.

TerraFormer uses policy-guided verifier feedback for dataset construction and model fine-tuning~\cite{terraformer2026}. This is a training-and-evaluation contribution that CloudCanvas does not implement. IaCGen instead emphasizes iterative formatting, syntax, and actual deployment feedback over 153 CloudFormation scenarios~\cite{iacgen2026}. Its results depend on the number of attempts, underlying models, and the definition of success. Neither benchmark is directly comparable to classification of synthetic graph fixtures. The preliminary cloud-agent study compares SDK, command-line, Terraform, and web agents on Azure VM tasks, then motivates synchronized state, policy enforcement, and human supervision~\cite{agents2025}. Differences in underlying models and task scope limit interpretation of its operational comparisons.

\subsection{Comparison and Research Gap}
Table~\ref{tab:related} compares six representative papers on the dimensions most relevant to this study. It records the type of evidence each contributes, rather than treating absent discussion of a feature as proof that the feature is unsupported. The remaining four papers supply complementary context on executable topology, runtime models, deployment abstractions, and agent supervision.
\begin{table}[t]
\caption{Representative related work and its relevance to CloudCanvas.}
\label{tab:related}
\centering\footnotesize
\begin{tabularx}{\linewidth}{p{.17\linewidth}XXX}
\toprule
Work and problem & Approach & Evaluation and result & Limitation and relevance \\
\midrule
Winery~\cite{winery2013}: topology design & Typed visual editor and repository & Moodle on two EC2 VMs demonstrates composition & Establishes canvas prior art with no controlled outcome comparison \\
DOML~\cite{doml2024}: model correctness & Model layers and Z3 property checks & Modified WordPress case exposed mapping and network errors & Limited generality across cases, but stronger model invariants \\
IaC-Eval~\cite{iaceval2024}: generated intent & Terraform plan with a policy oracle & 458 scenarios. Historical GPT-4 pass@1 of 19.36\% & AWS/Terraform scope motivates separate intent checks \\
Multi-IaC-Eval~\cite{multi2025}: template edits & Static and semantic checks & 337 CloudFormation, 171 Terraform, and 191 CDK templates & No live deployment test. Motivates preservation checks for edits \\
TerraFormer~\cite{terraformer2026}: generation quality & Policy feedback and fine-tuning & IaC-Eval correctness rose from 15.50\% to 31.44\% for the reported 14B condition & Depends on policy coverage. Motivates actionable diagnostics \\
IaCGen~\cite{iacgen2026}: deployability & Iterative cloud feedback & 153 scenarios. First-attempt deployability of 20.8--30.2\% across six models & Provider and model dependence calls for separate lifecycle measures \\
\bottomrule
\end{tabularx}
\end{table}
Commercial prior art strengthens this qualification. Massdriver documents visual composition of IaC bundles, JSON Schema input/output contracts, dependency checks, and managed deployment pipelines~\cite{massdriver}. Accordingly, visual composition plus schemas plus deployment is not, by itself, an original contribution. The specific contribution here is a source-grounded account and controlled conformance evaluation of CloudCanvas's shared intermediate graph and reviewed execution workflow. Its comparison baselines remove validation layers within that same representation. This isolates a tractable question while leaving usability, AI generation accuracy, enterprise governance, and live cloud reliability as separate evaluation targets.


The studies also differ in how they establish the expected result. Winery and OpenTOSCA demonstrate representative application workflows. CAMEL collects practitioner judgments. DOML checks declared properties on application models. IaC-Eval and TerraFormer depend on requirements encoded as policies, and IaCGen adds feedback from attempted deployment. Each method creates visibility into a different part of the system. A practitioner rating cannot explain which graph errors an implementation detects. A formal check of a particular property cannot establish that a user asked for that property. A successful deployment cannot establish that the resulting configuration reflects every requested restriction. These differences inform both the structure of the present experiment and its limitations.

Table~\ref{tab:related} therefore gives the reported benchmark values only with their original denominators and comparison conditions. For example, TerraFormer's gain refers to its tested model and verifier pipeline, and IaCGen's first-attempt rates describe actual deployment attempts for a different artifact format. Neither number provides a target for the percentage of labeled structural fixtures classified correctly here. A future CloudCanvas study could adapt the underlying evaluation techniques by writing independent requirement predicates, making review decisions observable, and checking actual resource outcomes against those predicates. Such a study would address a different research question from the current layer-by-layer conformance experiment.
